% GENERATED FILE. Edit canonical lessons, then run npm run generate:books.
% canonical-source-hash: fnv1a64:d8cacacdbb52c970
% canonical-lessons: TA-C530-kannamucci, TA-C530-pamparam, TA-C530-koli, TA-C530-keram, TA-C530-karrati

\chapter{Children's Games}
\label{ch:ta-children-s-games}
\hlchaptermodality{530}

\begin{chapteropening}
\textbf{By the end of this chapter:} \emph{I can name games children play outdoors.}

Complete the last lesson of chapter 530: I can name games children play outdoors.
\end{chapteropening}

\section[\texorpdfstring{\ta{கண்ணாமூச்சி} (kaṇṇāmūcci)}{கண்ணாமூச்சி (kaṇṇāmūcci)}]{\ta{கண்ணாமூச்சி} (kaṇṇāmūcci) — hide and seek}

\label{lesson:TA-C530-kannamucci}



\begin{quote}
Before the new one: say the Tamil for a die, then the Tamil for a seed-and-pit board game.
\end{quote}

\subsection*{You'll want to know: \ta{கண்ணாமூச்சி}}
\textbf{\ta{கண்ணாமூச்சி}} — \emph{kaṇṇāmūcci} — \textquotedblleft{}hide and seek\textquotedblright{}.

\begin{quote}
\ta{குழந்தைகள்} \ta{தோட்டத்தில்} \ta{கண்ணாமூச்சி} \ta{விளையாடினார்கள்}.
\end{quote}

\begin{quote}
\emph{kuḻantaikaḷ tōṭṭattil kaṇṇāmūcci viḷaiyāṭiṉārkaḷ.}
\end{quote}

\begin{quote}
The children played hide and seek in the garden.
\end{quote}

The seeker covers the eyes while the others hide.

\begin{cousinweb}
\textbf{\ta{கண்ணாமூச்சி}} (\emph{kaṇṇāmūcci}) begins with the native \ta{கண்} (\emph{kaṇ}), \textquotedblleft{}an eye\textquotedblright{} (DEDR 1159), since the seeker shuts the eyes; the formation of the rest is not settled.
\end{cousinweb}

\begin{grammarlens}[title={What you've built}]
One more word for describing what you lived and felt.
\end{grammarlens}

\subsection*{Your turn}
\begin{itemize}
\raggedright
  \item \emph{Say it:} \emph{kaṇṇāmūcci}
  \item \emph{Say it:} \emph{kaṇṇāmūcci}, once more
  \item \emph{Say it:} \emph{kuḻantaikaḷ tōṭṭattil kaṇṇāmūcci viḷaiyāṭiṉārkaḷ.} — The children played hide and seek in the garden.
  \item \emph{Recall:} say the Tamil for a die, then the Tamil for a seed-and-pit board game, then say \emph{kaṇṇāmūcci} again
\end{itemize}

\begin{tcolorbox}[breakable,colback=teal!4,colframe=teal!35!black,title={Before you move on}]
What does \ta{கண்ணாமூச்சி} mean? (\textquotedblleft{}Hide and seek\textquotedblright{}.) Say it once more.
\end{tcolorbox}
\section[\texorpdfstring{\ta{பம்பரம்} (pamparam)}{பம்பரம் (pamparam)}]{\ta{பம்பரம்} (pamparam) — a spinning top}

\label{lesson:TA-C530-pamparam}



\begin{quote}
Before the new one: say the Tamil for a seed-and-pit board game, then the Tamil for hide and seek.
\end{quote}

\subsection*{You'll want to know: \ta{பம்பரம்}}
\textbf{\ta{பம்பரம்}} — \emph{pamparam} — \textquotedblleft{}a spinning top\textquotedblright{}.

\begin{quote}
\ta{என்} \ta{தம்பி} \ta{தரையில்} \ta{பம்பரம்} \ta{சுற்றினான்}.
\end{quote}

\begin{quote}
\emph{eṉ tampi taraiyil pamparam cuṟṟiṉāṉ.}
\end{quote}

\begin{quote}
My little brother spun a top on the floor.
\end{quote}

One spins a top with \ta{சுற்று} (\emph{cuṟṟu}) or \ta{விடு}: \ta{பம்பரம்} \ta{விடு}.

\begin{cousinweb}
\textbf{\ta{பம்பரம்}} (\emph{pamparam}) is listed in DEDR 3930 as the Tamil word for a spinning top; the article marks a possible source in Sanskrit \emph{bhramaraka}, \textquotedblleft{}a humming top\textquotedblright{}, with a question mark, so the origin is not settled.
\end{cousinweb}

\begin{grammarlens}[title={What you've built}]
One more word for describing what you lived and felt.
\end{grammarlens}

\subsection*{Your turn}
\begin{itemize}
\raggedright
  \item \emph{Say it:} \emph{pamparam}
  \item \emph{Say it:} \emph{pamparam}, once more
  \item \emph{Say it:} \emph{eṉ tampi taraiyil pamparam cuṟṟiṉāṉ.} — My little brother spun a top on the floor.
  \item \emph{Recall:} say the Tamil for a seed-and-pit board game, then the Tamil for hide and seek, then say \emph{pamparam} again
\end{itemize}

\begin{tcolorbox}[breakable,colback=teal!4,colframe=teal!35!black,title={Before you move on}]
What does \ta{பம்பரம்} mean? (\textquotedblleft{}A spinning top\textquotedblright{}.) Say it once more.
\end{tcolorbox}
\section[\texorpdfstring{\ta{கோலி} (kōli)}{கோலி (kōli)}]{\ta{கோலி} (kōli) — marbles}

\label{lesson:TA-C530-koli}



\begin{quote}
Before the new one: say the Tamil for hide and seek, then the Tamil for a spinning top.
\end{quote}

\subsection*{You'll want to know: \ta{கோலி}}
\textbf{\ta{கோலி}} — \emph{kōli} — \textquotedblleft{}marbles\textquotedblright{}.

\begin{quote}
\ta{சிறுவர்கள்} \ta{மண்ணில்} \ta{கோலி} \ta{விளையாடினார்கள்}.
\end{quote}

\begin{quote}
\emph{ciṟuvarkaḷ maṇṇil kōli viḷaiyāṭiṉārkaḷ.}
\end{quote}

\begin{quote}
The boys played marbles in the dirt.
\end{quote}

One marble is a \ta{கோலிக்குண்டு}.

\begin{cousinweb}
The origin of \textbf{\ta{கோலி}} (\emph{kōli}) is not settled: no dictionary we can check gives it.
\end{cousinweb}

\begin{grammarlens}[title={What you've built}]
One more word for describing what you lived and felt.
\end{grammarlens}

\subsection*{Your turn}
\begin{itemize}
\raggedright
  \item \emph{Say it:} \emph{kōli}
  \item \emph{Say it:} \emph{kōli}, once more
  \item \emph{Say it:} \emph{ciṟuvarkaḷ maṇṇil kōli viḷaiyāṭiṉārkaḷ.} — The boys played marbles in the dirt.
  \item \emph{Recall:} say the Tamil for hide and seek, then the Tamil for a spinning top, then say \emph{kōli} again
\end{itemize}

\begin{tcolorbox}[breakable,colback=teal!4,colframe=teal!35!black,title={Before you move on}]
What does \ta{கோலி} mean? (\textquotedblleft{}Marbles\textquotedblright{}.) Say it once more.
\end{tcolorbox}
\section[\texorpdfstring{\ta{கேரம்} (kēram)}{கேரம் (kēram)}]{\ta{கேரம்} (kēram) — carrom}

\label{lesson:TA-C530-keram}



\begin{quote}
Before the new one: say the Tamil for a spinning top, then the Tamil for marbles.
\end{quote}

\subsection*{You'll want to know: \ta{கேரம்}}
\textbf{\ta{கேரம்}} — \emph{kēram} — \textquotedblleft{}carrom\textquotedblright{}.

\begin{quote}
\ta{மழைக்காலத்தில்} \ta{நாங்கள்} \ta{வீட்டில்} \ta{கேரம்} \ta{விளையாடுவோம்}.
\end{quote}

\begin{quote}
\emph{maḻaikkālattil nāṅkaḷ vīṭṭil kēram viḷaiyāṭuvōm.}
\end{quote}

\begin{quote}
In the rainy season we play carrom at home.
\end{quote}

A board game: players flick a striker to knock pieces into the four corner pockets.

\begin{cousinweb}
The origin of \textbf{\ta{கேரம்}} (\emph{kēram}) is not settled: it is the same word as English \emph{carrom}, but which way the name travelled is unclear.
\end{cousinweb}

\begin{grammarlens}[title={What you've built}]
One more word for describing what you lived and felt.
\end{grammarlens}

\subsection*{Your turn}
\begin{itemize}
\raggedright
  \item \emph{Say it:} \emph{kēram}
  \item \emph{Say it:} \emph{kēram}, once more
  \item \emph{Say it:} \emph{maḻaikkālattil nāṅkaḷ vīṭṭil kēram viḷaiyāṭuvōm.} — In the rainy season we play carrom at home.
  \item \emph{Recall:} say the Tamil for a spinning top, then the Tamil for marbles, then say \emph{kēram} again
\end{itemize}

\begin{tcolorbox}[breakable,colback=teal!4,colframe=teal!35!black,title={Before you move on}]
What does \ta{கேரம்} mean? (\textquotedblleft{}Carrom\textquotedblright{}.) Say it once more.
\end{tcolorbox}
\section[\texorpdfstring{\ta{காற்றாடி} (kāṟṟāṭi)}{காற்றாடி (kāṟṟāṭi)}]{\ta{காற்றாடி} (kāṟṟāṭi) — a pinwheel; a paper windmill}

\label{lesson:TA-C530-karrati}



\begin{quote}
Before the new one: say the Tamil for marbles, then the Tamil for carrom.
\end{quote}

\subsection*{You'll want to know: \ta{காற்றாடி}}
\textbf{\ta{காற்றாடி}} — \emph{kāṟṟāṭi} — \textquotedblleft{}a pinwheel; a paper windmill\textquotedblright{}.

\begin{quote}
\ta{திருவிழாவில்} \ta{அப்பா} \ta{எனக்கு} \ta{ஒரு} \ta{காற்றாடி} \ta{வாங்கினார்}.
\end{quote}

\begin{quote}
\emph{tiruviḻāvil appā eṉakku oru kāṟṟāṭi vāṅkiṉār.}
\end{quote}

\begin{quote}
At the festival Father bought me a pinwheel.
\end{quote}

Word for word \textquotedblleft{}wind-dancer\textquotedblright{}: \ta{காற்று} (\emph{kāṟṟu}) + \ta{ஆடி}, \textquotedblleft{}that which swings\textquotedblright{}.

\begin{cousinweb}
\textbf{\ta{காற்றாடி}} (\emph{kāṟṟāṭi}) is a native compound: \ta{காற்று} (\emph{kāṟṟu}), \textquotedblleft{}wind\textquotedblright{} (DEDR 1481), and \ta{ஆடி} (\emph{āṭi}), \textquotedblleft{}that which moves or swings\textquotedblright{}, from \ta{ஆடு}, \textquotedblleft{}to move, to dance\textquotedblright{}.
\end{cousinweb}

\begin{grammarlens}[title={What you've built}]
One more word for describing what you lived and felt.
\end{grammarlens}

\subsection*{Your turn}
\begin{itemize}
\raggedright
  \item \emph{Say it:} \emph{kāṟṟāṭi}
  \item \emph{Say it:} \emph{kāṟṟāṭi}, once more
  \item \emph{Say it:} \emph{tiruviḻāvil appā eṉakku oru kāṟṟāṭi vāṅkiṉār.} — At the festival Father bought me a pinwheel.
  \item \emph{Recall:} say the Tamil for marbles, then the Tamil for carrom, then say \emph{kāṟṟāṭi} again
\end{itemize}

\begin{tcolorbox}[breakable,colback=teal!4,colframe=teal!35!black,title={Before you move on}]
What does \ta{காற்றாடி} mean? (\textquotedblleft{}A pinwheel; a paper windmill\textquotedblright{}.) Say it once more.
\end{tcolorbox}
